%!TEX root=../Vorlage_DA.tex
%	########################################################
% 			INFO: Zitieren, Abbildungen, Listing
%	########################################################


%	--------------------------------------------------------
% 	Überschrift, Inhaltsverzeichnis
%	--------------------------------------------------------
\chapter*{INFO: Zitieren, Abbildungen, Listing}

In diesem Kapitel sind Beispiele angeführt, wie Abbildungen, Zitate und Quelltext zu verwenden sind.

Die Zitate und die Abbildungen werden automatisch in das Quellen- und Abbildungsverzeichnis übernommen. Das Quellenverzeichnis und das Abbildungsverzeichnis sind am Ende der Arbeit zu finden.   

\textbf{Beispiel: Icons--Beschreibung}
\setlength{\extrarowheight}{0.1cm}

\begin{tabularx}{\textwidth}{p{1.2cm} | X }

\vspace{-0.4cm}\includegraphics[width=1cm]{./media/images/htl_c_cmyk_rein.pdf} & Hier kommt dann die Beschreibung zum Icon dazu. Wenn es über mehrer Zeilen gehen sollte, dann erfolgt das mit dem \LaTeX--Befehl \inlinecode{newline}. 
\newline Hier geht die Beschreibung dann weiter \ldots
\\ 
\vspace{-0.4cm}\includegraphics[width=1cm]{./media/images/htl_c_cmyk_rein.pdf} & Hier kommt dann die Beschreibung zum Icon dazu. Wenn es über mehrer Zeilen gehen sollte, dann erfolgt das mit dem \LaTeX--Befehl \inlinecode{newline}. 
\newline Hier geht die Beschreibung dann weiter \ldots
\\
\vspace{-0.4cm}\includegraphics[width=1cm]{./media/images/htl_c_cmyk_rein.pdf} & Hier kommt dann die Beschreibung zum Icon dazu. Wenn es über mehrer Zeilen gehen sollte, dann erfolgt das mit dem \LaTeX--Befehl \inlinecode{newline}. 
\newline Hier geht die Beschreibung dann weiter \ldots
\\
\vspace{-0.4cm}\includegraphics[width=1cm]{./media/images/htl_c_cmyk_rein.pdf} & Hier kommt dann die Beschreibung zum Icon dazu. Wenn es über mehrer Zeilen gehen sollte, dann erfolgt das mit dem \LaTeX--Befehl \inlinecode{newline}. 
\newline Hier geht die Beschreibung dann weiter \ldots


\end{tabularx}


%	--------------------------------------------------------
% 	Inhalt 
%	--------------------------------------------------------
\section*{Abbildungen}


In der Abbildung \ref{htl01} auf Seite \pageref{htl01} ist das Logo der HTL im Farbmodus \emph{CMYK} zu sehen. Das Logo ist im \emph{.pdf}-Format gespeichert und setzt sich ausschließlich aus Pfaden zusammen.

%	--------------
%	Eine Abbildung
\begin{figure}[t]
	\centering
	\includegraphics[width=0.5\textwidth]{./media/images/htl_c_cmyk_rein.pdf}
  	\caption{Logo der HTL Braunau.}
  	\label{htl01}
\end{figure}


%	----------------
%	Zwei Abbildungen

\begin{figure}[bht]
  \centering
  \subfloat[Logo der HTL Braunau.]{\label{fig:a}\includegraphics[width=.4\textwidth]{./media/images/htl_c_cmyk_rein.pdf}}
  \qquad
  \subfloat[Logo der HTL Braunau.]{\label{fig:b}\includegraphics[width=.4\textwidth]{./media/images/htl_c_cmyk_rein.pdf}}
  \label{fig:canvas_01}
\caption[]{Zwei Logos der HTL Braunau mit in einer Abbildung zusammengefasst.}
\end{figure}


%	----------------
%	Drei Abbildungen
\begin{figure}[htb]
  \centering
  \subfloat[Logo der HTL Braunau.]{\label{fig:a}\includegraphics[width=.29\textwidth]{./media/images/htl_c_cmyk_rein.pdf}}
  \qquad
  \subfloat[Logo der HTL Braunau.]{\label{fig:b}\includegraphics[width=.29\textwidth]{./media/images/htl_c_cmyk_rein.pdf}}
\qquad
  \subfloat[Logo der HTL Braunau.]{\label{fig:b}\includegraphics[width=.29\textwidth]{./media/images/htl_c_cmyk_rein.pdf}}
  \label{fig:canvas_04}
\caption[]{Die Reihenfolge der Transformationen wirkt sich auf die Transformationsmatrix aus.}
\end{figure}


\section*{Zitate, Fußnoten, Literaturverzeichnis}

Die einzelnen Quellen sind in einer eigenen Datei hinterlegt. Konkret ist es bei diesem Dokument die Datei \href{"./includes/literatur.bib"}{\emph{literatur.bib}}, die sich im Verzeichnis \emph{includes} befindet.

Man sollte immer höflich bleiben ~\cite{autor:01}.\\

Wörtliche Zitate werden durch Anführungszeichen begonnen und beendet und kursiv geschrieben: \emph{\glqq We tend to think of navigating a website as clicking from page-to-page via some kind of global navigation that’s always visible. When it comes to a single page, we often think scrolling is the one and only way to move from one end to the next.\grqq}~\cite{smashing:01}.


Es ist zunehmend eine Kurzzitierweise in Fußnoten üblich: Nachname des Verfassers, Kurztitel, Seitenangabe.\footnote{John Doe, \glqq Verwendung von Fußnoten\grqq, Seite 12.}


%	%%%%%%%%%%%%%%%%%%%%%%%%%%%%%%%%%%%%%%%%%%%%%
%	Listings, Code
%	%%%%%%%%%%%%%%%%%%%%%%%%%%%%%%%%%%%%%%%%%%%%%
\section*{Listings, Code}
Wie aus dem Quelltext \ref{code01} auf Seite \pageref{code01} hervorgeht, \ldots


%	%%%%%%%%%%%%%%%%%%%%%%%%%%%%%%%%%%%%%%%%%%%%%
%	C++
\begin{figure}[H]
	\lstinputlisting[language=C++]{./media/code_snippets/bsp_01.cpp}
	\caption{Mein ersters C-Programm.}
  	\label{code01}
\end{figure}


\begin{lstlisting}[language=PHP]
<?php
if ($handle = opendir(realpath(__DIR__ . '/../userimages'))) {
    
    while (false !== ($file = readdir($handle))) {
        echo "$file//";
    }

    closedir($handle);
}
?>
\end{lstlisting}


\begin{lstlisting}[language=JavaScript]
// Calculates the mouse position.
function getMousePos(canvas, evt) {
	var rect = canvas.getBoundingClientRect();
	return {
		x: evt.clientX - rect.left,
		y: evt.clientY - rect.top
	};
}
\end{lstlisting}




%	%%%%%%%%%%%%%%%%%%%%%%%%%%%%%%%%%%%%%%%%%%%%%
%	CSS
\begin{figure}[H]
	\lstinputlisting[language=CSS]{./media/code_snippets/bsp.css}
	\caption{Kurzer CSS--Quelltext.}
  	\label{code02}
\end{figure}


%	%%%%%%%%%%%%%%%%%%%%%%%%%%%%%%%%%%%%%%%%%%%%%
%	HTML
\begin{figure}[H]
	\lstinputlisting[language=HTML]{./media/code_snippets/bsp.html}
	\caption{HTML--Quelltext}
  	\label{code03}
\end{figure}


%	%%%%%%%%%%%%%%%%%%%%%%%%%%%%%%%%%%%%%%%%%%%%%
%	HTML
\begin{figure}[H]
	\lstinputlisting[language=JavaScript]{./media/code_snippets/bsp.js}
	\caption{JavaScript--Quelltext}
  	\label{code04}
\end{figure}

%	%%%%%%%%%%%%%%%%%%%%%%%%%%%%%%%%%%%%%%%%%%%%%
%	Listings direkt im Programmfluss
Text vor einem Quelltext--Block \ldots

\begin{lstlisting}[language=CSS]
#idStyle.class {
	background-color: red;
	color: white;
}

.classStyle {
	color: blue;
}
\end{lstlisting}

\ldots Text nach einem Quelltext--Block.


\begin{lstlisting}[language=TypoScript]
10.
\end{lstlisting}

